\section{Simuladores Estocásticos}

Um simulador estocástico \( \mathcal{M}_S \) pode ser formalmente definido pelo seguinte mapeamento:
\begin{equation}
\mathcal{M}_S : \mathcal{D}_X \times \Omega \rightarrow \mathbb{R}, \qquad
(\mathbf{x}, \omega) \mapsto \mathcal{M}_S(\mathbf{x}, \omega),
\end{equation}
em que \( \mathbf{x} \in \mathcal{D}_X \subset \mathbb{R}^{M_X} \) é o vetor \( M_X \)-dimensional de parâmetros de entrada, e \( \omega \) é um evento aleatório em um espaço de probabilidade \( (\Omega, \mathcal{F}, \mathbb{P}) \), o qual representa a estocasticidade inerente do sistema.  

%Na prática, esta última decorre de variáveis aleatórias latentes \( \mathbf{Z}(\omega) \) (definidas sobre o mesmo espaço de probabilidade, com suporte \( \mathcal{D}_Z \subset \mathbb{R}^{M_Z} \)), que são parâmetros internos do simulador não explicitamente considerados como variáveis de entrada.  

Em outras palavras, um simulador estocástico é construído a partir de um modelo determinístico \( \mathcal{M} \) da seguinte forma:
\begin{equation}
\mathcal{M}_S(\mathbf{x}, \omega) = \mathcal{M}(\mathbf{x}, \mathbf{Z}(\omega)).
\label{eq3}
\end{equation}

Para um dado vetor de entrada \( \mathbf{x}_0 \), cada execução do simulador corresponde a uma realização diferente das variáveis latentes \( \mathbf{Z} \), digamos \( \mathbf{z}_0 \equiv \mathbf{Z}(\omega_0) \), e fornece uma realização particular da variável aleatória de saída \( Y_{\mathbf{x}_0} \equiv \mathcal{M}_S(\mathbf{x}_0, \cdot) \).  
Para caracterizar a distribuição desta variável de saída, devem ser geradas repetições, isto é, múltiplas execuções de \( \mathcal{M}_S \) com o mesmo vetor de entrada \( \mathbf{x}_0 \) e diferentes realizações \(\{\mathbf{z}_1, \ldots, \mathbf{z}_R\}\) das variáveis latentes.  Referir-nos-emos à função densidade de probabilidade de \( Y_{\mathbf{x}_0} \) como a \textit{distribuição condicional da resposta} dado \( \mathbf{x}_0 \).

 
\subsection{Análise de confiabilidade em simuladores estocásticos}

Seguindo o formalismo de \cite{Choe2015}, a probabilidade de falha associada a uma 
função limite estocástica \(g_s\) é definida como:
\begin{equation}
P_f = P \bigl( g_s(X, \omega) \le 0 \bigr) .
\end{equation}

Introduzindo a função determinística subjacente \(g\) e as variáveis latentes \(Z\), temos:
\begin{equation}
P_f = P \bigl( g(X, Z) \le 0 \bigr).
\label{eq5}
\end{equation}

Assumindo que \(X\) e \(Z\) são independentes e que o domínio de falha é definido no 
espaço conjunto \((X,Z)\),
\[
D_f = \{(x, z) : g(x, z) \le 0 \},
\]
podemos expandir a equação como:
\begin{equation}
P_f =
\int_{\mathbb{R}^{M_X}} \int_{\mathbb{R}^{M_Z}} \mathbf{1}_{D_f}(x, z) f_X(x) f_Z(z) \, dz \, dx ,
\label{eq6}
\end{equation}

%onde \(\mathbf{1}_{D_f}(x, z)\) é a função indicadora do domínio de falha, igual a 1 se 
%\(g(x, z) \le 0\) e 0 caso contrário. 

Esta definição integra a incerteza proveniente tanto dos parâmetros de entrada quanto das variáveis latentes.

Além disso, ao manipular a Eq. (\ref{eq6}), podemos derivar definições equivalentes de $P_f$ com base em duas perspectivas.  Para um dado \(\mathbf{x}_0\), a probabilidade condicional de falha é definida como:
\begin{equation}
P_Z \bigl( g(X, Z) \le 0 \mid X = \mathbf{x}_0 \bigr) =
\int_{\mathbb{R}^{M_Z}} \mathbf{1}_{D_f}(\mathbf{x}_0, z) f_Z(z) \, dz .
\end{equation}

 Para diferentes valores de \( x \), obtemos então uma \textit{função determinística de probabilidade condicional de falha}, que denotamos por \( s(x) \).

\begin{equation}
s(\mathbf{x}) = P_Z \bigl( g(X, Z) \le 0 \mid X = \mathbf{x} \bigr) ,
\label{eq8}
\end{equation}

Esta função reflete a probabilidade de falha devido à estocasticidade interna da própria função de estado limite.

Reformulando a Eq.(\ref{eq6}), podemos expressar \( P_f \) como uma função de \( s(x) \):

\begin{equation}
\begin{aligned}
P_f &= \int_{\mathbb{R}^{M_X}} \left( \int_{\mathbb{R}^{M_Z}} \mathbf{1}_{D_f}(x,z) f_Z(z) \, dz \right) f_X(x) \, dx \\
    &= \int_{\mathbb{R}^{M_X}} s(x) f_X(x) \, dx \\
    &\equiv \mathbb{E}_X \left[ s(X) \right] .
\end{aligned}
\label{eq9}
\end{equation}
%%%%%%%%%%%%%%
Isso nos permite decompor a estimação da probabilidade de falha em um processo de duas etapas, a saber: a caracterização da função \( s \) (por meio de equações semianalíticas) e o cálculo de seu valor esperado (ver detalhes na Seção~4).

Alternativamente, quando a estocasticidade latente é fixa, podemos expressar \( P_f \) como:
\begin{equation}
    P_f = \mathbb{E}_{\mathbf{Z}}\left[\mathbb{P}_{\mathbf{X}}\left(g(\mathbf{X},\mathbf{Z}) \leq 0 \mid \mathbf{Z} = z \right)\right].
    \label{eq10}
\end{equation}

O termo interno na Eq.~\eqref{eq10} representa a probabilidade de falha devido à incerteza no parâmetro de entrada para uma realização fixa da estocasticidade. Em outras palavras, \( P_f \) pode ser calculado como o valor esperado (sobre várias sementes aleatórias) das probabilidades de falha associadas a cada trajetória resultante.



\subsection{Expansão Estocástica em Caos Polinomial (SPCE)}
%\cite{Zhu2023_SPCE}}
A \emph{Stochastic Polynomial Chaos Expansion} (SPCE), introduzida por Zhu e Sudret \cite{Zhu2023_SPCE}, 
é  um emulador estocástico que permite modelar a distribuição de saída de $Y_x$.  
Nessa abordagem, uma variável latente artificial $Z$ é introduzida para representar a 
estocasticidade do simulador dentro do arcabouço das expansões em caos polinomial (PCE).

Essa formulação se baseia na \emph{transformada integral de probabilidade} (PIT), 
uma técnica estatística que permite transformar uma variável aleatória contínua em outra, 
possibilitando mapear uma variável conhecida $Z$ para a distribuição de saída desejada. 
Tipicamente, $Z$ é escolhida como uma variável com distribuição uniforme padrão ou normal padrão.

Seja $F_{Y|X}(y\,|\,x)$ a função de distribuição acumulada (CDF) associada a $Y_x$.  
A PIT permite expressar a relação entre $Y_x$ e a variável latente $Z$ como:
\begin{equation}
  Y_x \stackrel{d}{=} F_{Y|X}^{-1}\!\left(F_Z(Z)\,|\,x\right),
  \label{eq17}
\end{equation}
 
onde $F_Z$ é a CDF da variável latente $Z$, e o símbolo $\stackrel{d}{=}$ indica igualdade em distribuição.

A Eq. (\ref{eq17}) implica que $Y_x$ pode ser representada como uma função determinística tanto da entrada $x$ 
quanto da variável latente $Z$. Essa relação determinística é então modelada por uma expansão em caos 
polinomial no espaço $(X, Z)$:
\begin{equation}
  F_{Y|X}^{-1}\!\left(F_Z(Z)\,|\,X = x\right)
\stackrel{d}{=}
\sum_{\boldsymbol{\alpha} \in \mathbb{N}^{M_X + 1}} 
c_{\boldsymbol{\alpha}} \, \psi_{\boldsymbol{\alpha}}(x, Z).
\label{eq18}
\end{equation}
 
Para um dado $x$, a expansão acima é uma função da variável latente $Z$, tornando a resposta $Y_x$ 
uma variável aleatória. Considerando um esquema de truncamento $\mathcal{A}$, obtém-se:
\begin{equation}
  Y_x \stackrel{d}{\approx} 
\tilde{Y}_x = 
\sum_{\boldsymbol{\alpha} \in \mathcal{A}} 
c_{\boldsymbol{\alpha}} \, \psi_{\boldsymbol{\alpha}}(x, Z).
\label{eq19}
\end{equation}
 
A SPCE, portanto, fornece uma estrutura unificada para representar simuladores estocásticos, 
permitindo capturar tanto a variabilidade quanto a dependência funcional das saídas em relação às entradas. 
Mais detalhes teóricos e exemplos de aplicação podem ser encontrados em Zhu e Sudret~\cite{Zhu2023_SPCE}.

%\subsection{Regularização e formulação completa da SPCE \cite{Zhu2023_SPCE}}

O uso direto da Eq.~(\ref{eq19}) pode levar a problemas como singularidades nas funções densidade de probabilidade (PDFs).  
Para contornar essa limitação, Zhu e Sudret~\cite{Zhu2023_SPCE} introduziram um componente aditivo de ruído
\(\varepsilon \sim \mathcal{N}(0, \sigma^2)\), que atua, na prática, como um termo de regularização.  
O emulador SPCE resultante é então expresso como:
\begin{equation}
Y_x \stackrel{d}{\approx} 
\tilde{Y}_x =
\sum_{\boldsymbol{\alpha} \in \mathcal{A}}
c_{\boldsymbol{\alpha}} \, \psi_{\boldsymbol{\alpha}}(x, Z)
+ \varepsilon.
\label{eq:spce_noise}
\end{equation}

A PDF de \(\tilde{Y}_x\) resulta da convolução entre a PDF da expansão PCE e a PDF gaussiana do ruído.  
Assumindo que a variância do ruído \(\sigma^2\) é conhecida, a PDF correspondente pode ser escrita como:
%%%%%%%%%%%%%%%%%%%%%%%%%%%%Detalhamento da covolução
Seja o termo determinístico-estocástico da PCE definido por
\[
X \equiv 
\sum_{\boldsymbol{\alpha} \in \mathcal{A}}
c_{\boldsymbol{\alpha}} \, 
\psi_{\boldsymbol{\alpha}}(x, Z),
\]
e o ruído aditivo
\[
Y \equiv \varepsilon \sim \mathcal{N}(0,\sigma^{2}).
\]
Como \(X\) e \(Y\) são independentes, a densidade da soma
\(\tilde{Y}_x = X + Y\) é dada pela convolução das respectivas PDFs:
\[
f_{\tilde{Y}_x}(y)
=
\int_{-\infty}^{\infty}
f_X(t)\, f_Y(y - t)\, dt.
\]

No caso particular em que \(Y\) é gaussiana, sua PDF é
\[
f_Y(y) = \frac{1}{\sigma}\,\phi\!\left(\frac{y}{\sigma}\right),
\]
em que \(\phi(\cdot)\) é a densidade da normal padrão. Substituindo esta
expressão na convolução, obtém-se
\[
f_{\tilde{Y}_x}(y)
=
\int_{-\infty}^{\infty}
f_X(t)\,
\frac{1}{\sigma}
\phi\!\left(
\frac{y - t}{\sigma}
\right)
dt.
\]

Como a variável \(X\) depende do vetor aleatório \(Z\), sua densidade
pode ser escrita condicionalmente como
\[
f_X(t)
=
\int_{\mathcal{D}_Z}
\delta\!\left(
t -
\sum_{\boldsymbol{\alpha} \in \mathcal{A}}
c_{\boldsymbol{\alpha}} \psi_{\boldsymbol{\alpha}}(x, z)
\right)
f_Z(z)\, dz.
\]
Substituindo esta expressão na integral de convolução e integrando em
relação a \(t\), obtém-se finalmente
%\[
%f_{\tilde{Y}_x}(y)
%=
%\int_{\mathcal{D}_Z}
%\frac{1}{\sigma}
%\phi\!\left(
%\frac{
%y -\sum_{\boldsymbol{\alpha} \in \mathcal{A}}
%c_{\boldsymbol{\alpha}}
%\psi_{\boldsymbol{\alpha}}(x, z)}{\sigma}\right)
%f_Z(z)\, dz,
%\]

%%%%%%%%%%%%%%%%
\begin{equation}
f_{\tilde{Y}_x}(y) =
\int_{\mathcal{D}_Z}
\frac{1}{\sigma}
\phi\!\left(
\frac{
y -
\sum_{\boldsymbol{\alpha} \in \mathcal{A}} 
c_{\boldsymbol{\alpha}} \psi_{\boldsymbol{\alpha}}(x, z)
}{\sigma}
\right)
f_Z(z) \, dz,
\label{eq21}
\end{equation}
em que \(\phi(\cdot)\) representa a função densidade da normal padrão.
%%%%%%%%%%%%%%
\subsection{Cálculo de \(s(x)\) a partir do emulador SPCE}

Para o SPCE, a CDF associada a uma previsão \(\tilde{Y}_x\), denotada como \(F_{\tilde{Y}_x}(y)\), é obtida integrando a PDF introduzida na Eq. (\ref{eq21}). Como esta PDF é uma integral unidimensional, ela pode ser avaliada utilizando quadratura de Gauss (Hermite)da seguinte forma:
\begin{equation}
f_{\tilde{Y}_x}(y) = \int_{D_Z} \frac{1}{\sigma} \varphi \left( \frac{y - \sum_{\alpha \in A} c_\alpha \psi_\alpha(x,z)}{\sigma} \right) f_Z(z) dz
\approx \sum_{j=1}^{N_Q} \frac{w^{(j)}}{\sqrt{2\pi}\sigma} \exp \left( - \frac{\left(y - \sum_{\alpha \in A} c_\alpha \psi_\alpha \big(x, z^{(j)}\big)\right)^2}{2\sigma^2} \right),
\end{equation}
onde \(N_Q\) é o número de pontos de integração, \(z^{(j)}\) é o \(j\)-ésimo ponto de integração, e \(w^{(j)}\) é o peso correspondente, ambos associados à função de peso \(f_Z\), que na prática é a PDF uniforme ou Gaussiana, para a qual os pontos e pesos de integração são bem conhecidos.

A aproximação introduzida pelo esquema de integração numérica implica que \(f_{\tilde{Y}_x}(y)\) é aproximada por uma mistura de \(N_Q\) distribuições Gaussianas
\[
\mathcal{N} \Bigg( \sum_{\alpha \in A} c_\alpha \psi_\alpha \big(x, z^{(j)}\big), \sigma^2 \Bigg), \quad j=1,\dots,N_Q.
\]

Com os correspondentes pesos $w^{(j)}$. A CDF (função de distribuição acumulada) dessa aproximação é direta e pode ser escrita como:

\begin{align}
F_{Y_x}(y) 
&\approx \int_{-\infty}^{y} 
\sum_{j=1}^{N_Q} 
\frac{w^{(j)}}{\sqrt{2\pi}\sigma}
\exp\!\left[
-\frac{\left(
y - \sum_{\alpha \in \mathcal{A}} c_\alpha \psi_\alpha(x, z^{(j)})
\right)^2}{2\sigma^2}
\right] dy \nonumber \\[6pt]
&\approx 
\sum_{j=1}^{N_Q} 
w^{(j)} 
\int_{-\infty}^{y}
\frac{1}{\sqrt{2\pi}\sigma}
\exp\!\left[
-\frac{\left(
y - \sum_{\alpha \in \mathcal{A}} c_\alpha \psi_\alpha(x, z^{(j)})
\right)^2}{2\sigma^2}
\right] dy \nonumber \\[6pt]
&\approx 
\sum_{j=1}^{N_Q}
w^{(j)} 
\Phi\!\left(
\frac{
y - \sum_{\alpha \in \mathcal{A}} c_\alpha \psi_\alpha(x, z^{(j)})
}{\sigma}
\right),
\label{eq:29}
\end{align}

onde $\Phi$ é a CDF Gaussiana padrão.

Finalmente, a função de probabilidade condicional de falha é dada por $F_{Y_x}(0)$, ou seja:

\begin{equation}
s(x) \approx
\sum_{j=1}^{N_Q}
w^{(j)} 
\Phi\!\left(
-\frac{
\sum_{\alpha \in \mathcal{A}} c_\alpha \psi_\alpha(x, z^{(j)})
}{\sigma}
\right).
\label{eq:30}
\end{equation}

Uma vez que o modelo SPCE é treinado, o cálculo de 
$\sum_{\alpha \in \mathcal{A}} c_\alpha \psi_\alpha(x, z)$ 
é computacionalmente barato.  
A precisão dessa abordagem depende fortemente do número de pontos de quadratura utilizados.  
Neste trabalho, utilizou-se $N_Q = 100$, o que, com base em nossos testes, resulta em erro numérico desprezível.
 